
% ==============================================================
% Section 6: Threats to Validity
% Author: Quân (RW)
% File:   paper/sections/06_threats.tex
% ==============================================================

\section{Threats to Validity}\label{sec:threats}

\subsection{Internal Validity}

\textbf{Prompt change between phases.}
Phase~1 used prompt \texttt{v1.0} and Phase~2 used
\texttt{v3.0-candidate-s2r-relaxed}.  The change was motivated by the
systematic S2R labeling error identified in Phase~1 (code snippets
labeled \textit{Missing} by \texttt{v1.0} but \textit{Sufficient} by
annotators).  Because the two phases ran under different prompt
versions and Phase~2 has no human labels, the effect of \texttt{v3.0}
on agreement is untested.  All agreement results in
Section~\ref{sec:results} characterise \texttt{v1.0} only; claims
about \texttt{v3.0} are limited to operational feasibility.

\textbf{Limited annotator overlap.}
Annotator~2 independently labeled 10 of the 30 Phase~1 reports
(33\%).  Human--human Kappa is therefore based on this 10-report
overlap, not on all 30 issues.  While the result ($\kappa = 1.000$)
indicates consistent rubric application within that subset, it may not
reflect agreement on the full range of reports in the annotated set.

\textbf{OB Kappa undefined.}
Every report in the annotated subset received \textit{Sufficient} for
OB from both the model and the annotators.  When a single label
exhausts the distribution, chance-expected agreement equals 1 and
Cohen's Kappa is undefined ($0/0$).  We characterise OB by raw
agreement (100\%) and exclude it from averaged Kappa, as noted in
Sections~\ref{sec:method:metrics} and~\ref{sec:results}.  The
undefined coefficient should not be read as a positive finding; it
reflects a lack of label variation on this dimension in this sample.

\textbf{Unrecorded sampling seed.}
The random seed used to sample the 30 pilot issues was not recorded.
The sampled reports are committed to the repository so our specific
results are verifiable, but the exact sample cannot be regenerated by
an independent third party.

\subsection{External Validity}

\textbf{Dataset scope.}
The annotated subset comprises 30 issues from three repositories
(pandas, scikit-learn, VS~Code).  The distribution of OB, EB, and S2R
quality in these repositories may not represent the broader population
of GitHub bug reports, enterprise projects, or other issue-tracking
platforms.  Caution is warranted in generalizing the Kappa values
beyond the evaluated repositories.

\textbf{Single model and configuration.}
This study evaluates GPT-4o mini at temperature~$0.0$ only.  Results
may differ with other LLMs, higher temperatures, or alternative prompt
designs.  Generalizing the findings to other models or configurations
requires additional experimentation.

\subsection{Construct Validity}

\textbf{Rubric boundary between \textit{Ambiguous} and
\textit{Sufficient}.}
The boundary between these two labels---particularly for S2R---involves
a judgment about the completeness of a reproduction path.  The two
annotators applied it consistently in this sample, but annotators with
different domain backgrounds may draw the boundary differently.  The
rubric was not validated with annotators outside the project team.

\textbf{Cohen's Kappa as the sole agreement metric.}
Kappa accounts for chance agreement but does not reveal the direction
of disagreement.  The label distribution in
Table~\ref{tab:label_dist} provides directional evidence (the model
under-labels \textit{Missing} relative to annotators), but a full
per-label confusion matrix per dimension would give a more precise
picture of where and how the model errs---information that would
directly guide further prompt refinement.
